
\textbf{Related Work}
Temporal Graph Neural Networks (TGNNs) are categorized into continuous and discrete time dynamic graph methods (CTDG and DTDG). CTDG methods process event streams with timestamps using neighbor sampling to model evolving relations. They include TGAT \citep{tgat} with self-attention and temporal encoding, TGN \citep{tgn} with memory modules, DyGFormer \citep{dygformer} using multi-head attention on temporal patches, and CTAN \citep{ctan} employing ODEs. DTDG methods operate on regularly-spaced graph snapshots and typically use recurrent neural networks to track history; popular methods include GC-LSTM \citep{gclstm}, T-GCN \citep{tgcn}, and EvolveGCN \citep{evolvegcn}). 
Recently, \cite{huangutg} compared CTDG and DTDG methods and found that DTDG methods sacrifice accuracy for efficiency by computing on coarser-level snapshot information. Subsequently, they showed that DTDG methods can operate on CTDG datasets by using the Unified Temporal Graph (UTG) framework. In \Cref{sec:tgrab}, we evaluate and compare the above methods from both continuous and discrete-time approaches on \benchmarkname datasets to analyze their capabilities in capturing core temporal patterns. Dedicated benchmarks \citep{Jodie_Kumar_2019, uci, enron, social-evo} have significantly improved TGNN evaluation, though \cite{poursafaei2022betterevaluationdynamiclink} noted overly optimistic results due to simple negative edges and introduced new sampling strategies and diverse datasets. Subsequently, TGB \citep{tgb} and TGB 2.0 \citep{tgb2.0} presented larger, challenging datasets for link and node prediction, while TGB-Seq \citep{tgb-seq} proposed real-world  with complex sequential dynamics and low edge repeatability, revealing limitations in current TGNN generalization. Our work complements these recent efforts by constructing synthetic tasks to more exactly pinpoint the current functional weaknesses of TGNNs.


%  \emph{Temporal Graph Neural Networks}~(TGNNs) can be broadly categorized based on whether they operate on continuous-time or discrete-time representations of temporal graphs \citep{kazemi2020representation}. Continuous-Time Dynamic Graph (CTDG) performs computation on a stream of events, each of which is associated with a timestamp. They typically employ neighbor sampling to have access to historical information and subsequently model evolving relations. Many recent architectures operate on this principle: for instance, TGAT \citep{tgat} leverages self-attention with temporal encoding to weigh the importance of historical interactions; TGN \citep{tgn} generalizes this with memory modules to capture long-range temporal dependencies; DyGFormer \citep{dygformer} computes node embeddings by applying multi-head attention \citep{transformer} to temporal patches of a node's first-hop interactions; and CTAN \citep{ctan} frames temporal graph learning as an Ordinary Differential Equation (ODE) to model long-range dependencies.



% \looseness-1 Discrete-Time Dynamic Graph (DTDG) methods, conversely, operate on a sequence of graph snapshots sampled at regularly-spaced time intervals \citep{kazemi2020representation}. As opposed to the neighbour sampling done by CTDG methods, DTDG methods typically use a recurrent neural network to keep track of historical information.  Early TGNN approaches such as GC-LSTM \citep{gclstm}, T-GCN \citep{tgcn}, and EvolveGCN \citep{evolvegcn} exemplify this, often combining static GNNs to process the structure within each snapshot with recurrent neural networks to model the evolution across snapshots. Recently, \cite{huangutg} compared CTDG and DTDG methods and found that DTDG methods sacrifice accuracy for efficiency by computing on coarser-level snapshot information. Subsequently, they showed that DTDG methods can operate on CTDG datasets by using the Unified Temporal Graph (UTG) framework. In \Cref{sec:tgrab}, we evaluate and compare the above methods from both continuous and discrete-time approaches on \benchmarkname datasets to analyze their capabilities in capturing core temporal patterns.

% \looseness-1 \textbf{Benchmarks for Temporal Graph Learning.} The rigorous evaluation of TGNNs has been significantly advanced by the use of dedicated benchmarks \citep{Jodie_Kumar_2019, uci, enron, social-evo}. Following this, \cite{poursafaei2022betterevaluationdynamiclink} identified that results on previous benchmarks are overly optimistic due to easy negative edges and proposed novel negative sampling strategies in addition to contributing a range of datasets that aim to increase the domain diversity of current benchmarks. Shortly after, TGB \citep{tgb} and TGB 2.0 \citep{tgb2.0} introduced significantly larger datasets for link prediction and node-level tasks that proved to be challenging for current TGNNs. TGB-Seq \citep{tgb-seq} proposed a collection of real-world datasets that featured complex sequential dynamics and low edge repeatability that demanded higher generalization power from models. However, they again show that current TGNNs achieve disappointing results. Our work complements these recent efforts by constructing synthetic tasks to more exactly pinpoint the current functional weaknesses of TGNNs.
